% Options for packages loaded elsewhere
\PassOptionsToPackage{unicode}{hyperref}
\PassOptionsToPackage{hyphens}{url}
\PassOptionsToPackage{dvipsnames,svgnames,x11names}{xcolor}
%
\documentclass[
  11pt,
]{article}
\usepackage{amsmath,amssymb}
\usepackage{lmodern}
\usepackage{iftex}
\ifPDFTeX
  \usepackage[T1]{fontenc}
  \usepackage[utf8]{inputenc}
  \usepackage{textcomp} % provide euro and other symbols
\else % if luatex or xetex
  \usepackage{unicode-math}
  \defaultfontfeatures{Scale=MatchLowercase}
  \defaultfontfeatures[\rmfamily]{Ligatures=TeX,Scale=1}
\fi
% Use upquote if available, for straight quotes in verbatim environments
\IfFileExists{upquote.sty}{\usepackage{upquote}}{}
\IfFileExists{microtype.sty}{% use microtype if available
  \usepackage[]{microtype}
  \UseMicrotypeSet[protrusion]{basicmath} % disable protrusion for tt fonts
}{}
\makeatletter
\@ifundefined{KOMAClassName}{% if non-KOMA class
  \IfFileExists{parskip.sty}{%
    \usepackage{parskip}
  }{% else
    \setlength{\parindent}{0pt}
    \setlength{\parskip}{6pt plus 2pt minus 1pt}}
}{% if KOMA class
  \KOMAoptions{parskip=half}}
\makeatother
\usepackage{xcolor}
\IfFileExists{xurl.sty}{\usepackage{xurl}}{} % add URL line breaks if available
\IfFileExists{bookmark.sty}{\usepackage{bookmark}}{\usepackage{hyperref}}
\hypersetup{
  pdftitle={Project Proposal (Week 2), Energy-Aware Ventilation \& Climate Control},
  pdfauthor={Example team: A. Andersson \& B. Berg},
  colorlinks=true,
  linkcolor={RoyalBlue},
  filecolor={Maroon},
  citecolor={Blue},
  urlcolor={Blue},
  pdfcreator={LaTeX via pandoc}}
\urlstyle{same} % disable monospaced font for URLs
\usepackage[a4paper,margin=2.4cm]{geometry}
\usepackage{longtable,booktabs,array}
\usepackage{calc} % for calculating minipage widths
% Correct order of tables after \paragraph or \subparagraph
\usepackage{etoolbox}
\makeatletter
\patchcmd\longtable{\par}{\if@noskipsec\mbox{}\fi\par}{}{}
\makeatother
% Allow footnotes in longtable head/foot
\IfFileExists{footnotehyper.sty}{\usepackage{footnotehyper}}{\usepackage{footnote}}
\makesavenoteenv{longtable}
\setlength{\emergencystretch}{3em} % prevent overfull lines
\providecommand{\tightlist}{%
  \setlength{\itemsep}{0pt}\setlength{\parskip}{0pt}}
\setcounter{secnumdepth}{-\maxdimen} % remove section numbering
\usepackage{newunicodechar}
\usepackage{textcomp}
\usepackage{amssymb}
\newunicodechar{°}{\textdegree}
\newunicodechar{₂}{\textsubscript{2}}
\newunicodechar{→}{\ensuremath{\rightarrow}}
\newunicodechar{≤}{\ensuremath{\leq}}
\newunicodechar{≥}{\ensuremath{\geq}}
\newunicodechar{≈}{\ensuremath{\approx}}
\newunicodechar{✓}{\ensuremath{\checkmark}}
\newunicodechar{–}{--}
\newunicodechar{—}{---}
\usepackage{etoolbox}
\AtBeginEnvironment{longtable}{\footnotesize}
\AtBeginEnvironment{tabular}{\footnotesize}
\usepackage{fancyhdr}
\pagestyle{fancy}\fancyhf{}
\rhead{\small D7065E\,--\,Project Proposal}
\lhead{\small Example: Week 2 Proposal}
\cfoot{\thepage}
\renewcommand{\headrulewidth}{0.4pt}
\usepackage{titlesec}
\titleformat{\section}{\Large\bfseries}{\thesection}{0.6em}{}
\titleformat{\subsection}{\large\bfseries}{\thesubsection}{0.5em}{}
\usepackage{microtype}
\setlength{\emergencystretch}{2em}
\ifLuaTeX
  \usepackage{selnolig}  % disable illegal ligatures
\fi

\title{Project Proposal (Week 2), Energy-Aware Ventilation \& Climate
Control}
\usepackage{etoolbox}
\makeatletter
\providecommand{\subtitle}[1]{% add subtitle to \maketitle
  \apptocmd{\@title}{\par {\large #1 \par}}{}{}
}
\makeatother
\subtitle{D7065E, worked example}
\author{Example team: A. Andersson \& B. Berg}
\date{}

\begin{document}
\maketitle

\begin{quote}
\textbf{About this document.} A short, worked example of the
\textbf{Week 2 proposal}. It is deliberately provisional, its job is
to get a \emph{direction} approved, not to lock down the design. Some
early choices here (REST transport, an ML occupancy model) are revised
later in the final report; that is expected and good.
\end{quote}

\hypertarget{summary}{%
\subsection{Summary}\label{summary}}

\begin{longtable}[]{@{}
  >{\raggedright\arraybackslash}p{(\columnwidth - 2\tabcolsep) * \real{0.5000}}
  >{\raggedright\arraybackslash}p{(\columnwidth - 2\tabcolsep) * \real{0.5000}}@{}}
\toprule
\endhead
\textbf{Team} & A. Andersson, B. Berg \\
\textbf{Use case} & Energy-aware ventilation \& climate control \\
\textbf{Pitch} & Ventilate and heat rooms based on who is actually
present, to keep air quality and comfort high while cutting energy. \\
\textbf{Repository} &
\texttt{github.com/example-team/d7065e-ventilation} \\
\bottomrule
\end{longtable}

\hypertarget{use-case-and-why}{%
\subsection{1. Use case, and why}\label{use-case-and-why}}

This project controls ventilation and heating in a BuildSim office
floor. It was chosen because it has a genuine tension, air quality and
comfort pull one way, energy the other, which makes the autonomous
decision interesting, and because BuildSim provides rooms and occupancy
to work with.

\hypertarget{scenario}{%
\subsection{2. Scenario}\label{scenario}}

A meeting fills a room; CO₂ and temperature rise; the system detects it
and increases ventilation, holding CO₂ below 1000 ppm; when the room
empties it falls back to a low-energy setback. The same loop runs in
every room.

\hypertarget{what-well-build-rough-sketch}{%
\subsection{\texorpdfstring{3. What will be built \emph{(rough
sketch)}}{3. What will be built (rough sketch)}}\label{what-well-build-rough-sketch}}

\begin{itemize}
\tightlist
\item
  \textbf{Sensors:} temperature, CO₂, and occupancy per room, registered
  with BuildSim.
\item
  \textbf{Autonomous decision-maker:} sets ventilation level and heating
  setpoint per room.
\item
  \textbf{Actuators:} HVAC setpoint and ventilation damper, writing
  state back to BuildSim.
\item
  \textbf{Data pipeline / storage:} collect sensor readings, store
  recent history, and provide data to the autonomous service and the
  dashboard.
\item
  \textbf{Dashboard:} overlay state and decisions on the BuildSim 3D
  view.
\end{itemize}

Each major component (sensors, autonomous service, actuators, storage,
and dashboard) runs as a \textbf{separate process} communicating through
\textbf{defined interfaces}, and talks to BuildSim through its API.

\hypertarget{simulating-the-building}{%
\subsection{4. Simulating the building}\label{simulating-the-building}}

BuildSim provides the rooms; the physics and the people are simulated as
follows.

\begin{itemize}
\tightlist
\item
  \textbf{Physical processes:} a simple per-room model, temperature
  from a heat balance (envelope loss, body heat, HVAC), and CO₂ that
  rises with the number of people and falls with ventilation. Actuator
  commands feed back into the model so the loop closes.
\item
  \textbf{People movement (occupancy):} sample a head-count per room
  conditioned on \textbf{room type and time of day}, offices steady
  through the day, the fika room busy around 11:30, conference rooms full
  during meetings, empty at night. (A richer option that could be added
  later: move individual people between rooms along BuildSim's
  \textbf{navigation graph}.)
\end{itemize}

The models are intentionally simple, but \textbf{actuator actions must
influence future sensor readings} so that the control loop closes.

\hypertarget{the-autonomous-part}{%
\subsection{5. The autonomous part}\label{the-autonomous-part}}

It decides each room's fan level and setpoint from CO₂, temperature, and
occupancy. \textbf{Starting approach:} a small ML model to predict
occupancy, feeding a simple controller. Plain rules may be used instead
if the data turns out too thin to train on.

\hypertarget{biggest-unknowns-risks}{%
\subsection{6. Biggest unknowns \& risks}\label{biggest-unknowns-risks}}

\begin{itemize}
\tightlist
\item
  Whether two weeks of simulated data is enough to train the occupancy
  model.
\item
  Whether sending readings to the controller over REST keeps up as
  rooms are added.
\item
  Whether the room physics is realistic enough for the decisions to mean
  anything.
\end{itemize}

\hypertarget{first-step-rough-plan}{%
\subsection{7. First step \& rough plan}\label{first-step-rough-plan}}

The first step is \textbf{one room working end-to-end}, one set of
sensors → controller → actuator → BuildSim, to prove the loop closes.
After that, more rooms, storage, and the dashboard are added, and failure
handling and evaluation follow toward the Week 8 demo.

\end{document}
